% =====================================================================
%  Anatomía de un LLM — Fase 0 — Módulo 0.2
%  Conjuntos, funciones y relaciones
%  Compilar: latexmk -r ../latexmkrc P02-conjuntos-funciones.tex
% =====================================================================
\documentclass[11pt,a4paper,twoside,openany]{book}
\usepackage{../estilo-apuntes}

\newcommand{\ApuntesCurso}{Anatomía de un LLM}
\newcommand{\ApuntesModulo}{Módulo 0.2}

% --- Notación propia del módulo -------------------------------------------
\newcommand{\N}{\mathbb{N}}
\newcommand{\Z}{\mathbb{Z}}
\newcommand{\Q}{\mathbb{Q}}
\newcommand{\Pa}{\mathcal{P}}
\newcommand{\id}{\mathrm{id}}
\newcommand{\cad}[1]{\texttt{#1}}
\newcommand{\vacia}{\varepsilon}
\newcommand{\dsim}{\mathbin{\triangle}}

\hypersetup{pdftitle={Módulo 0.2. Conjuntos, funciones y relaciones (resumen)},pdfauthor={Anatomía de un LLM}}

\makeatletter
\newlabel{ej:t1-pertenencia}{{1.2}{1}{}{ejemplo.2}{}}
\newlabel{ej:t1-pertenencia@cref}{{[ejemplo][2][]1.2}{[1][1][]1}}
\newlabel{fig:t1-venn}{{1.1}{3}{Diagramas de Venn de la intersección y de la primera ley de De Morgan. Un diagrama ilustra una igualdad, pero no la demuestra}{figure.1.1}{}}
\newlabel{fig:t1-venn@cref}{{[figure][1][1]1.1}{[1][3][]3}}
\newlabel{ej:t1-doble}{{1.10}{3}{Una igualdad por doble inclusión}{ejemplo.10}{}}
\newlabel{ej:t1-doble@cref}{{[ejemplo][10][]1.10}{[1][3][]3}}
\newlabel{eje:1-pert}{{1.1}{4}{}{ejercicio.1.1}{}}
\newlabel{eje:1-pert@cref}{{[ejercicio][1][1]1.1}{[1][4][]4}}
\newlabel{eje:1-card}{{1.2}{4}{}{ejercicio.1.2}{}}
\newlabel{eje:1-card@cref}{{[ejercicio][2][1]1.2}{[1][4][]4}}
\newlabel{eje:1-extension}{{1.3}{4}{}{ejercicio.1.3}{}}
\newlabel{eje:1-extension@cref}{{[ejercicio][3][1]1.3}{[1][4][]4}}
\newlabel{eje:1-partes}{{1.4}{4}{}{ejercicio.1.4}{}}
\newlabel{eje:1-partes@cref}{{[ejercicio][4][1]1.4}{[1][4][]4}}
\newlabel{eje:1-cartesiano}{{1.5}{4}{}{ejercicio.1.5}{}}
\newlabel{eje:1-cartesiano@cref}{{[ejercicio][5][1]1.5}{[1][4][]4}}
\newlabel{eje:1-demorgan-num}{{1.6}{4}{}{ejercicio.1.6}{}}
\newlabel{eje:1-demorgan-num@cref}{{[ejercicio][6][1]1.6}{[1][4][]4}}
\newlabel{eje:1-vf}{{1.7}{5}{}{ejercicio.1.7}{}}
\newlabel{eje:1-vf@cref}{{[ejercicio][7][1]1.7}{[1][4][]5}}
\newlabel{eje:1-dif-union}{{1.8}{5}{}{ejercicio.1.8}{}}
\newlabel{eje:1-dif-union@cref}{{[ejercicio][8][1]1.8}{[1][5][]5}}
\newlabel{eje:1-union-inclusion}{{1.9}{5}{}{ejercicio.1.9}{}}
\newlabel{eje:1-union-inclusion@cref}{{[ejercicio][9][1]1.9}{[1][5][]5}}
\newlabel{eje:1-producto-inter}{{1.10}{5}{}{ejercicio.1.10}{}}
\newlabel{eje:1-producto-inter@cref}{{[ejercicio][10][1]1.10}{[1][5][]5}}
\newlabel{eje:1-simetrica}{{1.11}{5}{}{ejercicio.1.11}{}}
\newlabel{eje:1-simetrica@cref}{{[ejercicio][11][1]1.11}{[1][5][]5}}
\newlabel{eje:1-partes-inter}{{1.12}{5}{}{ejercicio.1.12}{}}
\newlabel{eje:1-partes-inter@cref}{{[ejercicio][12][1]1.12}{[1][5][]5}}
\newlabel{eje:1-familias}{{1.13}{5}{}{ejercicio.1.13}{}}
\newlabel{eje:1-familias@cref}{{[ejercicio][13][1]1.13}{[1][5][]5}}
\newlabel{eje:1-cadenas}{{1.14}{5}{}{ejercicio.1.14}{}}
\newlabel{eje:1-cadenas@cref}{{[ejercicio][14][1]1.14}{[1][5][]5}}
\newlabel{ej:t2-funciones}{{2.2}{6}{}{ejemplo.2}{}}
\newlabel{ej:t2-funciones@cref}{{[ejemplo][2][]2.2}{[1][6][]6}}
\newlabel{ej:t2-biendef}{{2.3}{6}{Funciones bien definidas}{ejemplo.3}{}}
\newlabel{ej:t2-biendef@cref}{{[ejemplo][3][]2.3}{[1][6][]6}}
\newlabel{ej:t2-cuadrado}{{2.8}{7}{}{ejemplo.8}{}}
\newlabel{ej:t2-cuadrado@cref}{{[ejemplo][8][]2.8}{[1][7][]7}}
\newlabel{eje:2-esfuncion}{{2.1}{8}{}{ejercicio.2.1}{}}
\newlabel{eje:2-esfuncion@cref}{{[ejercicio][1][2]2.1}{[1][8][]8}}
\newlabel{eje:2-tabla}{{2.2}{8}{}{ejercicio.2.2}{}}
\newlabel{eje:2-tabla@cref}{{[ejercicio][2][2]2.2}{[1][8][]8}}
\newlabel{eje:2-contar}{{2.3}{8}{}{ejercicio.2.3}{}}
\newlabel{eje:2-contar@cref}{{[ejercicio][3][2]2.3}{[1][8][]8}}
\newlabel{eje:2-cuadrado}{{2.4}{8}{}{ejercicio.2.4}{}}
\newlabel{eje:2-cuadrado@cref}{{[ejercicio][4][2]2.4}{[1][8][]8}}
\newlabel{eje:2-suelo}{{2.5}{8}{}{ejercicio.2.5}{}}
\newlabel{eje:2-suelo@cref}{{[ejercicio][5][2]2.5}{[1][8][]8}}
\newlabel{eje:2-longitud}{{2.6}{9}{}{ejercicio.2.6}{}}
\newlabel{eje:2-longitud@cref}{{[ejercicio][6][2]2.6}{[1][9][]9}}
\newlabel{eje:2-recorrido}{{2.7}{9}{}{ejercicio.2.7}{}}
\newlabel{eje:2-recorrido@cref}{{[ejercicio][7][2]2.7}{[1][9][]9}}
\newlabel{eje:2-preim-ui}{{2.8}{9}{}{ejercicio.2.8}{}}
\newlabel{eje:2-preim-ui@cref}{{[ejercicio][8][2]2.8}{[1][9][]9}}
\newlabel{eje:2-imagen-ui}{{2.9}{9}{}{ejercicio.2.9}{}}
\newlabel{eje:2-imagen-ui@cref}{{[ejercicio][9][2]2.9}{[1][9][]9}}
\newlabel{eje:2-ida-vuelta}{{2.10}{9}{}{ejercicio.2.10}{}}
\newlabel{eje:2-ida-vuelta@cref}{{[ejercicio][10][2]2.10}{[1][9][]9}}
\newlabel{eje:2-preim-comp}{{2.11}{9}{}{ejercicio.2.11}{}}
\newlabel{eje:2-preim-comp@cref}{{[ejercicio][11][2]2.11}{[1][9][]9}}
\newlabel{eje:2-caracteristica}{{2.12}{9}{}{ejercicio.2.12}{}}
\newlabel{eje:2-caracteristica@cref}{{[ejercicio][12][2]2.12}{[1][9][]9}}
\newlabel{eje:2-biendef}{{2.13}{9}{}{ejercicio.2.13}{}}
\newlabel{eje:2-biendef@cref}{{[ejercicio][13][2]2.13}{[1][9][]9}}
\newlabel{eje:2-argmax}{{2.14}{9}{}{ejercicio.2.14}{}}
\newlabel{eje:2-argmax@cref}{{[ejercicio][14][2]2.14}{[1][9][]9}}
\newlabel{fig:t3-flechas}{{3.1}{10}{Diagramas de flechas de funciones entre conjuntos finitos}{figure.3.1}{}}
\newlabel{fig:t3-flechas@cref}{{[figure][1][3]3.1}{[1][10][]10}}
\newlabel{ej:t3-2n1}{{3.2}{11}{}{ejemplo.2}{}}
\newlabel{ej:t3-2n1@cref}{{[ejemplo][2][]3.2}{[1][11][]11}}
\newlabel{ej:t3-cuadrado}{{3.3}{11}{Dominio y codominio importan}{ejemplo.3}{}}
\newlabel{ej:t3-cuadrado@cref}{{[ejemplo][3][]3.3}{[1][11][]11}}
\newlabel{ej:t3-suma}{{3.4}{11}{Una función de dos variables}{ejemplo.4}{}}
\newlabel{ej:t3-suma@cref}{{[ejemplo][4][]3.4}{[1][11][]11}}
\newlabel{ej:t3-cadenas}{{3.5}{11}{Funciones sobre cadenas}{ejemplo.5}{}}
\newlabel{ej:t3-cadenas@cref}{{[ejemplo][5][]3.5}{[1][11][]11}}
\newlabel{ej:t3-sigmoide}{{3.8}{11}{La sigmoide}{ejemplo.8}{}}
\newlabel{ej:t3-sigmoide@cref}{{[ejemplo][8][]3.8}{[1][11][]11}}
\newlabel{eje:3-tablas}{{3.1}{12}{}{ejercicio.3.1}{}}
\newlabel{eje:3-tablas@cref}{{[ejercicio][1][3]3.1}{[1][12][]12}}
\newlabel{eje:3-error}{{3.2}{13}{}{ejercicio.3.2}{}}
\newlabel{eje:3-error@cref}{{[ejercicio][2][3]3.2}{[1][12][]13}}
\newlabel{eje:3-clasificar}{{3.3}{13}{}{ejercicio.3.3}{}}
\newlabel{eje:3-clasificar@cref}{{[ejercicio][3][3]3.3}{[1][13][]13}}
\newlabel{eje:3-5n2}{{3.4}{13}{}{ejercicio.3.4}{}}
\newlabel{eje:3-5n2@cref}{{[ejercicio][4][3]3.4}{[1][13][]13}}
\newlabel{eje:3-homografica}{{3.5}{13}{}{ejercicio.3.5}{}}
\newlabel{eje:3-homografica@cref}{{[ejercicio][5][3]3.5}{[1][13][]13}}
\newlabel{eje:3-lineal-z2}{{3.6}{13}{}{ejercicio.3.6}{}}
\newlabel{eje:3-lineal-z2@cref}{{[ejercicio][6][3]3.6}{[1][13][]13}}
\newlabel{eje:3-2m3n}{{3.7}{13}{}{ejercicio.3.7}{}}
\newlabel{eje:3-2m3n@cref}{{[ejercicio][7][3]3.7}{[1][13][]13}}
\newlabel{eje:3-NZ}{{3.8}{13}{}{ejercicio.3.8}{}}
\newlabel{eje:3-NZ@cref}{{[ejercicio][8][3]3.8}{[1][13][]13}}
\newlabel{eje:3-creciente}{{3.9}{13}{}{ejercicio.3.9}{}}
\newlabel{eje:3-creciente@cref}{{[ejercicio][9][3]3.9}{[1][13][]13}}
\newlabel{eje:3-cadenas}{{3.10}{13}{}{ejercicio.3.10}{}}
\newlabel{eje:3-cadenas@cref}{{[ejercicio][10][3]3.10}{[1][13][]13}}
\newlabel{eje:3-codigos}{{3.11}{13}{}{ejercicio.3.11}{}}
\newlabel{eje:3-codigos@cref}{{[ejercicio][11][3]3.11}{[1][13][]13}}
\newlabel{eje:3-contar}{{3.12}{14}{}{ejercicio.3.12}{}}
\newlabel{eje:3-contar@cref}{{[ejercicio][12][3]3.12}{[1][14][]14}}
\newlabel{eje:3-finitos}{{3.13}{14}{}{ejercicio.3.13}{}}
\newlabel{eje:3-finitos@cref}{{[ejercicio][13][3]3.13}{[1][14][]14}}
\newlabel{eje:3-softsign}{{3.14}{14}{}{ejercicio.3.14}{}}
\newlabel{eje:3-softsign@cref}{{[ejercicio][14][3]3.14}{[1][14][]14}}
\newlabel{eje:3-tanh}{{3.15}{14}{}{ejercicio.3.15}{}}
\newlabel{eje:3-tanh@cref}{{[ejercicio][15][3]3.15}{[1][14][]14}}
\newlabel{eje:3-fibras}{{3.16}{14}{}{ejercicio.3.16}{}}
\newlabel{eje:3-fibras@cref}{{[ejercicio][16][3]3.16}{[1][14][]14}}
\newlabel{ej:t4-comp}{{4.2}{15}{}{ejemplo.2}{}}
\newlabel{ej:t4-comp@cref}{{[ejemplo][2][]4.2}{[1][15][]15}}
\newlabel{ej:t4-inversa}{{4.10}{16}{}{ejemplo.10}{}}
\newlabel{ej:t4-inversa@cref}{{[ejemplo][10][]4.10}{[1][16][]16}}
\newlabel{eje:4-calc-comp}{{4.1}{17}{}{ejercicio.4.1}{}}
\newlabel{eje:4-calc-comp@cref}{{[ejercicio][1][4]4.1}{[1][17][]17}}
\newlabel{eje:4-vf}{{4.2}{17}{}{ejercicio.4.2}{}}
\newlabel{eje:4-vf@cref}{{[ejercicio][2][4]4.2}{[1][17][]17}}
\newlabel{eje:4-tablas}{{4.3}{17}{}{ejercicio.4.3}{}}
\newlabel{eje:4-tablas@cref}{{[ejercicio][3][4]4.3}{[1][17][]17}}
\newlabel{eje:4-inversas-reales}{{4.4}{17}{}{ejercicio.4.4}{}}
\newlabel{eje:4-inversas-reales@cref}{{[ejercicio][4][4]4.4}{[1][17][]17}}
\newlabel{eje:4-inversa-z2}{{4.5}{17}{}{ejercicio.4.5}{}}
\newlabel{eje:4-inversa-z2@cref}{{[ejercicio][5][4]4.5}{[1][17][]17}}
\newlabel{eje:4-comp-sobre}{{4.6}{17}{}{ejercicio.4.6}{}}
\newlabel{eje:4-comp-sobre@cref}{{[ejercicio][6][4]4.6}{[1][17][]17}}
\newlabel{eje:4-cambridge}{{4.7}{17}{}{ejercicio.4.7}{}}
\newlabel{eje:4-cambridge@cref}{{[ejercicio][7][4]4.7}{[1][17][]17}}
\newlabel{eje:4-iteradas}{{4.8}{18}{}{ejercicio.4.8}{}}
\newlabel{eje:4-iteradas@cref}{{[ejercicio][8][4]4.8}{[1][18][]18}}
\newlabel{eje:4-inv-comp}{{4.9}{18}{}{ejercicio.4.9}{}}
\newlabel{eje:4-inv-comp@cref}{{[ejercicio][9][4]4.9}{[1][18][]18}}
\newlabel{eje:4-derecha}{{4.10}{18}{}{ejercicio.4.10}{}}
\newlabel{eje:4-derecha@cref}{{[ejercicio][10][4]4.10}{[1][18][]18}}
\newlabel{eje:4-involucion}{{4.11}{18}{}{ejercicio.4.11}{}}
\newlabel{eje:4-involucion@cref}{{[ejercicio][11][4]4.11}{[1][18][]18}}
\newlabel{eje:4-cancelar}{{4.12}{18}{}{ejercicio.4.12}{}}
\newlabel{eje:4-cancelar@cref}{{[ejercicio][12][4]4.12}{[1][18][]18}}
\newlabel{eje:4-capas}{{4.13}{18}{}{ejercicio.4.13}{}}
\newlabel{eje:4-capas@cref}{{[ejercicio][13][4]4.13}{[1][18][]18}}
\newlabel{eje:4-tokenizador}{{4.14}{18}{}{ejercicio.4.14}{}}
\newlabel{eje:4-tokenizador@cref}{{[ejercicio][14][4]4.14}{[1][18][]18}}
\newlabel{ej:t5-modn}{{5.7}{20}{Enteros módulo $n$}{ejemplo.7}{}}
\newlabel{ej:t5-modn@cref}{{[ejemplo][7][]5.7}{[1][20][]20}}
\newlabel{ej:t5-ordenes}{{5.10}{20}{}{ejemplo.10}{}}
\newlabel{ej:t5-ordenes@cref}{{[ejemplo][10][]5.10}{[1][20][]20}}
\newlabel{fig:t5-hasse}{{5.1}{21}{Diagrama de Hasse de $(\Pa (\{a,b,c\}),\subseteq )$. El mínimo es $\emptyset $ y el máximo, $\{a,b,c\}$}{figure.5.1}{}}
\newlabel{fig:t5-hasse@cref}{{[figure][1][5]5.1}{[1][21][]21}}
\newlabel{eje:5-propiedades}{{5.1}{21}{}{ejercicio.5.1}{}}
\newlabel{eje:5-propiedades@cref}{{[ejercicio][1][5]5.1}{[1][21][]21}}
\newlabel{eje:5-enteros}{{5.2}{21}{}{ejercicio.5.2}{}}
\newlabel{eje:5-enteros@cref}{{[ejercicio][2][5]5.2}{[1][21][]21}}
\newlabel{eje:5-contar}{{5.3}{22}{}{ejercicio.5.3}{}}
\newlabel{eje:5-contar@cref}{{[ejercicio][3][5]5.3}{[1][21][]22}}
\newlabel{eje:5-clases}{{5.4}{22}{}{ejercicio.5.4}{}}
\newlabel{eje:5-clases@cref}{{[ejercicio][4][5]5.4}{[1][22][]22}}
\newlabel{eje:5-particion}{{5.5}{22}{}{ejercicio.5.5}{}}
\newlabel{eje:5-particion@cref}{{[ejercicio][5][5]5.5}{[1][22][]22}}
\newlabel{eje:5-bell}{{5.6}{22}{}{ejercicio.5.6}{}}
\newlabel{eje:5-bell@cref}{{[ejercicio][6][5]5.6}{[1][22][]22}}
\newlabel{eje:5-modn}{{5.7}{22}{}{ejercicio.5.7}{}}
\newlabel{eje:5-modn@cref}{{[ejercicio][7][5]5.7}{[1][22][]22}}
\newlabel{eje:5-nucleo}{{5.8}{22}{}{ejercicio.5.8}{}}
\newlabel{eje:5-nucleo@cref}{{[ejercicio][8][5]5.8}{[1][22][]22}}
\newlabel{eje:5-divide}{{5.9}{22}{}{ejercicio.5.9}{}}
\newlabel{eje:5-divide@cref}{{[ejercicio][9][5]5.9}{[1][22][]22}}
\newlabel{eje:5-mas3}{{5.10}{22}{}{ejercicio.5.10}{}}
\newlabel{eje:5-mas3@cref}{{[ejercicio][10][5]5.10}{[1][22][]22}}
\newlabel{eje:5-fracciones}{{5.11}{22}{}{ejercicio.5.11}{}}
\newlabel{eje:5-fracciones@cref}{{[ejercicio][11][5]5.11}{[1][22][]22}}
\newlabel{eje:5-orden-partes}{{5.12}{22}{}{ejercicio.5.12}{}}
\newlabel{eje:5-orden-partes@cref}{{[ejercicio][12][5]5.12}{[1][22][]22}}
\newlabel{eje:5-hasse}{{5.13}{22}{}{ejercicio.5.13}{}}
\newlabel{eje:5-hasse@cref}{{[ejercicio][13][5]5.13}{[1][22][]22}}
\newlabel{eje:5-lex}{{5.14}{23}{}{ejercicio.5.14}{}}
\newlabel{eje:5-lex@cref}{{[ejercicio][14][5]5.14}{[1][22][]23}}
\newlabel{eje:5-error}{{5.15}{23}{}{ejercicio.5.15}{}}
\newlabel{eje:5-error@cref}{{[ejercicio][15][5]5.15}{[1][23][]23}}
\newlabel{ej:t6-pares}{{6.2}{24}{Un conjunto infinito y una parte propia}{ejemplo.2}{}}
\newlabel{ej:t6-pares@cref}{{[ejemplo][2][]6.2}{[1][24][]24}}
\newlabel{ej:t6-restos}{{6.4}{25}{Restos}{ejemplo.4}{}}
\newlabel{ej:t6-restos@cref}{{[ejemplo][4][]6.4}{[1][25][]25}}
\newlabel{fig:t6-cantor}{{6.1}{26}{Recorrido de $\N \times \N $ por diagonales. Cada círculo muestra la posición $\pi (m,n)$ del par $(m,n)$}{figure.6.1}{}}
\newlabel{fig:t6-cantor@cref}{{[figure][1][6]6.1}{[1][25][]26}}
\newlabel{eje:6-vf}{{6.1}{27}{}{ejercicio.6.1}{}}
\newlabel{eje:6-vf@cref}{{[ejercicio][1][6]6.1}{[1][27][]27}}
\newlabel{eje:6-textos}{{6.2}{27}{}{ejercicio.6.2}{}}
\newlabel{eje:6-textos@cref}{{[ejercicio][2][6]6.2}{[1][27][]27}}
\newlabel{eje:6-meses}{{6.3}{27}{}{ejercicio.6.3}{}}
\newlabel{eje:6-meses@cref}{{[ejercicio][3][6]6.3}{[1][27][]27}}
\newlabel{eje:6-bits}{{6.4}{27}{}{ejercicio.6.4}{}}
\newlabel{eje:6-bits@cref}{{[ejercicio][4][6]6.4}{[1][27][]27}}
\newlabel{eje:6-restos}{{6.5}{27}{}{ejercicio.6.5}{}}
\newlabel{eje:6-restos@cref}{{[ejercicio][5][6]6.5}{[1][27][]27}}
\newlabel{eje:6-puntos}{{6.6}{27}{}{ejercicio.6.6}{}}
\newlabel{eje:6-puntos@cref}{{[ejercicio][6][6]6.6}{[1][27][]27}}
\newlabel{eje:6-compresion}{{6.7}{27}{}{ejercicio.6.7}{}}
\newlabel{eje:6-compresion@cref}{{[ejercicio][7][6]6.7}{[1][27][]27}}
\newlabel{eje:6-pi}{{6.8}{28}{}{ejercicio.6.8}{}}
\newlabel{eje:6-pi@cref}{{[ejercicio][8][6]6.8}{[1][27][]28}}
\newlabel{eje:6-shortlex}{{6.9}{28}{}{ejercicio.6.9}{}}
\newlabel{eje:6-shortlex@cref}{{[ejercicio][9][6]6.9}{[1][28][]28}}
\newlabel{eje:6-diagonal}{{6.10}{28}{}{ejercicio.6.10}{}}
\newlabel{eje:6-diagonal@cref}{{[ejercicio][10][6]6.10}{[1][28][]28}}
\newlabel{eje:6-subconjN}{{6.11}{28}{}{ejercicio.6.11}{}}
\newlabel{eje:6-subconjN@cref}{{[ejercicio][11][6]6.11}{[1][28][]28}}
\newlabel{eje:6-producto}{{6.12}{28}{}{ejercicio.6.12}{}}
\newlabel{eje:6-producto@cref}{{[ejercicio][12][6]6.12}{[1][28][]28}}
\newlabel{eje:6-Q}{{6.13}{28}{}{ejercicio.6.13}{}}
\newlabel{eje:6-Q@cref}{{[ejercicio][13][6]6.13}{[1][28][]28}}
\newlabel{eje:6-PA}{{6.14}{28}{}{ejercicio.6.14}{}}
\newlabel{eje:6-PA@cref}{{[ejercicio][14][6]6.14}{[1][28][]28}}
\newlabel{eje:6-intervalos}{{6.15}{28}{}{ejercicio.6.15}{}}
\newlabel{eje:6-intervalos@cref}{{[ejercicio][15][6]6.15}{[1][28][]28}}
\newlabel{ap:soluciones}{{A}{29}{Soluciones de los ejercicios}{appendix.A}{}}
\newlabel{ap:soluciones@cref}{{[appendix][1][2147483647]A}{[1][29][]29}}
\newlabel{sol:eje:1-pert}{{A.1}{29}{Tema 1}{section*.4}{}}
\newlabel{sol:eje:1-pert@cref}{{[subappendix][1][2147483647,1]A.1}{[1][29][]29}}
\newlabel{sol:eje:1-card}{{A.1}{29}{Tema 1}{section*.5}{}}
\newlabel{sol:eje:1-card@cref}{{[subappendix][1][2147483647,1]A.1}{[1][29][]29}}
\newlabel{sol:eje:1-extension}{{A.1}{30}{Tema 1}{section*.6}{}}
\newlabel{sol:eje:1-extension@cref}{{[subappendix][1][2147483647,1]A.1}{[1][30][]30}}
\newlabel{sol:eje:1-partes}{{A.1}{30}{Tema 1}{section*.7}{}}
\newlabel{sol:eje:1-partes@cref}{{[subappendix][1][2147483647,1]A.1}{[1][30][]30}}
\newlabel{sol:eje:1-cartesiano}{{A.1}{30}{Tema 1}{section*.8}{}}
\newlabel{sol:eje:1-cartesiano@cref}{{[subappendix][1][2147483647,1]A.1}{[1][30][]30}}
\newlabel{sol:eje:1-demorgan-num}{{A.1}{31}{Tema 1}{section*.9}{}}
\newlabel{sol:eje:1-demorgan-num@cref}{{[subappendix][1][2147483647,1]A.1}{[1][31][]31}}
\newlabel{sol:eje:1-vf}{{A.1}{31}{Tema 1}{section*.10}{}}
\newlabel{sol:eje:1-vf@cref}{{[subappendix][1][2147483647,1]A.1}{[1][31][]31}}
\newlabel{sol:eje:1-dif-union}{{A.1}{31}{Tema 1}{section*.11}{}}
\newlabel{sol:eje:1-dif-union@cref}{{[subappendix][1][2147483647,1]A.1}{[1][31][]31}}
\newlabel{sol:eje:1-union-inclusion}{{A.1}{32}{Tema 1}{section*.12}{}}
\newlabel{sol:eje:1-union-inclusion@cref}{{[subappendix][1][2147483647,1]A.1}{[1][32][]32}}
\newlabel{sol:eje:1-producto-inter}{{A.1}{32}{Tema 1}{section*.13}{}}
\newlabel{sol:eje:1-producto-inter@cref}{{[subappendix][1][2147483647,1]A.1}{[1][32][]32}}
\newlabel{sol:eje:1-simetrica}{{A.1}{32}{Tema 1}{section*.14}{}}
\newlabel{sol:eje:1-simetrica@cref}{{[subappendix][1][2147483647,1]A.1}{[1][32][]32}}
\newlabel{sol:eje:1-partes-inter}{{A.1}{33}{Tema 1}{section*.15}{}}
\newlabel{sol:eje:1-partes-inter@cref}{{[subappendix][1][2147483647,1]A.1}{[1][33][]33}}
\newlabel{sol:eje:1-familias}{{A.1}{33}{Tema 1}{section*.16}{}}
\newlabel{sol:eje:1-familias@cref}{{[subappendix][1][2147483647,1]A.1}{[1][33][]33}}
\newlabel{sol:eje:1-cadenas}{{A.1}{34}{Tema 1}{section*.17}{}}
\newlabel{sol:eje:1-cadenas@cref}{{[subappendix][1][2147483647,1]A.1}{[1][34][]34}}
\newlabel{sol:eje:2-esfuncion}{{A.2}{34}{Tema 2}{section*.18}{}}
\newlabel{sol:eje:2-esfuncion@cref}{{[subappendix][2][2147483647,1]A.2}{[1][34][]34}}
\newlabel{sol:eje:2-tabla}{{A.2}{34}{Tema 2}{section*.19}{}}
\newlabel{sol:eje:2-tabla@cref}{{[subappendix][2][2147483647,1]A.2}{[1][34][]34}}
\newlabel{sol:eje:2-contar}{{A.2}{35}{Tema 2}{section*.20}{}}
\newlabel{sol:eje:2-contar@cref}{{[subappendix][2][2147483647,1]A.2}{[1][35][]35}}
\newlabel{sol:eje:2-cuadrado}{{A.2}{35}{Tema 2}{section*.21}{}}
\newlabel{sol:eje:2-cuadrado@cref}{{[subappendix][2][2147483647,1]A.2}{[1][35][]35}}
\newlabel{sol:eje:2-suelo}{{A.2}{35}{Tema 2}{section*.22}{}}
\newlabel{sol:eje:2-suelo@cref}{{[subappendix][2][2147483647,1]A.2}{[1][35][]35}}
\newlabel{sol:eje:2-longitud}{{A.2}{36}{Tema 2}{section*.23}{}}
\newlabel{sol:eje:2-longitud@cref}{{[subappendix][2][2147483647,1]A.2}{[1][36][]36}}
\newlabel{sol:eje:2-recorrido}{{A.2}{36}{Tema 2}{section*.24}{}}
\newlabel{sol:eje:2-recorrido@cref}{{[subappendix][2][2147483647,1]A.2}{[1][36][]36}}
\newlabel{sol:eje:2-preim-ui}{{A.2}{36}{Tema 2}{section*.25}{}}
\newlabel{sol:eje:2-preim-ui@cref}{{[subappendix][2][2147483647,1]A.2}{[1][36][]36}}
\newlabel{sol:eje:2-imagen-ui}{{A.2}{36}{Tema 2}{section*.26}{}}
\newlabel{sol:eje:2-imagen-ui@cref}{{[subappendix][2][2147483647,1]A.2}{[1][36][]36}}
\newlabel{sol:eje:2-ida-vuelta}{{A.2}{37}{Tema 2}{section*.27}{}}
\newlabel{sol:eje:2-ida-vuelta@cref}{{[subappendix][2][2147483647,1]A.2}{[1][37][]37}}
\newlabel{sol:eje:2-preim-comp}{{A.2}{37}{Tema 2}{section*.28}{}}
\newlabel{sol:eje:2-preim-comp@cref}{{[subappendix][2][2147483647,1]A.2}{[1][37][]37}}
\newlabel{sol:eje:2-caracteristica}{{A.2}{37}{Tema 2}{section*.29}{}}
\newlabel{sol:eje:2-caracteristica@cref}{{[subappendix][2][2147483647,1]A.2}{[1][37][]37}}
\newlabel{sol:eje:2-biendef}{{A.2}{38}{Tema 2}{section*.30}{}}
\newlabel{sol:eje:2-biendef@cref}{{[subappendix][2][2147483647,1]A.2}{[1][38][]38}}
\newlabel{sol:eje:2-argmax}{{A.2}{38}{Tema 2}{section*.31}{}}
\newlabel{sol:eje:2-argmax@cref}{{[subappendix][2][2147483647,1]A.2}{[1][38][]38}}
\newlabel{sol:eje:3-tablas}{{A.3}{39}{Tema 3}{section*.32}{}}
\newlabel{sol:eje:3-tablas@cref}{{[subappendix][3][2147483647,1]A.3}{[1][38][]39}}
\newlabel{sol:eje:3-error}{{A.3}{39}{Tema 3}{section*.33}{}}
\newlabel{sol:eje:3-error@cref}{{[subappendix][3][2147483647,1]A.3}{[1][39][]39}}
\newlabel{sol:eje:3-clasificar}{{A.3}{39}{Tema 3}{section*.34}{}}
\newlabel{sol:eje:3-clasificar@cref}{{[subappendix][3][2147483647,1]A.3}{[1][39][]39}}
\newlabel{sol:eje:3-5n2}{{A.3}{40}{Tema 3}{section*.35}{}}
\newlabel{sol:eje:3-5n2@cref}{{[subappendix][3][2147483647,1]A.3}{[1][40][]40}}
\newlabel{sol:eje:3-homografica}{{A.3}{40}{Tema 3}{section*.36}{}}
\newlabel{sol:eje:3-homografica@cref}{{[subappendix][3][2147483647,1]A.3}{[1][40][]40}}
\newlabel{sol:eje:3-lineal-z2}{{A.3}{40}{Tema 3}{section*.37}{}}
\newlabel{sol:eje:3-lineal-z2@cref}{{[subappendix][3][2147483647,1]A.3}{[1][40][]40}}
\newlabel{sol:eje:3-2m3n}{{A.3}{41}{Tema 3}{section*.38}{}}
\newlabel{sol:eje:3-2m3n@cref}{{[subappendix][3][2147483647,1]A.3}{[1][41][]41}}
\newlabel{sol:eje:3-NZ}{{A.3}{41}{Tema 3}{section*.39}{}}
\newlabel{sol:eje:3-NZ@cref}{{[subappendix][3][2147483647,1]A.3}{[1][41][]41}}
\newlabel{sol:eje:3-creciente}{{A.3}{41}{Tema 3}{section*.40}{}}
\newlabel{sol:eje:3-creciente@cref}{{[subappendix][3][2147483647,1]A.3}{[1][41][]41}}
\newlabel{sol:eje:3-cadenas}{{A.3}{42}{Tema 3}{section*.41}{}}
\newlabel{sol:eje:3-cadenas@cref}{{[subappendix][3][2147483647,1]A.3}{[1][42][]42}}
\newlabel{sol:eje:3-codigos}{{A.3}{42}{Tema 3}{section*.42}{}}
\newlabel{sol:eje:3-codigos@cref}{{[subappendix][3][2147483647,1]A.3}{[1][42][]42}}
\newlabel{sol:eje:3-contar}{{A.3}{42}{Tema 3}{section*.43}{}}
\newlabel{sol:eje:3-contar@cref}{{[subappendix][3][2147483647,1]A.3}{[1][42][]42}}
\newlabel{sol:eje:3-finitos}{{A.3}{43}{Tema 3}{section*.44}{}}
\newlabel{sol:eje:3-finitos@cref}{{[subappendix][3][2147483647,1]A.3}{[1][43][]43}}
\newlabel{sol:eje:3-softsign}{{A.3}{43}{Tema 3}{section*.45}{}}
\newlabel{sol:eje:3-softsign@cref}{{[subappendix][3][2147483647,1]A.3}{[1][43][]43}}
\newlabel{sol:eje:3-tanh}{{A.3}{44}{Tema 3}{section*.46}{}}
\newlabel{sol:eje:3-tanh@cref}{{[subappendix][3][2147483647,1]A.3}{[1][43][]44}}
\newlabel{sol:eje:3-fibras}{{A.3}{44}{Tema 3}{section*.47}{}}
\newlabel{sol:eje:3-fibras@cref}{{[subappendix][3][2147483647,1]A.3}{[1][44][]44}}
\newlabel{sol:eje:4-calc-comp}{{A.4}{44}{Tema 4}{section*.48}{}}
\newlabel{sol:eje:4-calc-comp@cref}{{[subappendix][4][2147483647,1]A.4}{[1][44][]44}}
\newlabel{sol:eje:4-vf}{{A.4}{45}{Tema 4}{section*.49}{}}
\newlabel{sol:eje:4-vf@cref}{{[subappendix][4][2147483647,1]A.4}{[1][45][]45}}
\newlabel{sol:eje:4-tablas}{{A.4}{45}{Tema 4}{section*.50}{}}
\newlabel{sol:eje:4-tablas@cref}{{[subappendix][4][2147483647,1]A.4}{[1][45][]45}}
\newlabel{sol:eje:4-inversas-reales}{{A.4}{45}{Tema 4}{section*.51}{}}
\newlabel{sol:eje:4-inversas-reales@cref}{{[subappendix][4][2147483647,1]A.4}{[1][45][]45}}
\newlabel{sol:eje:4-inversa-z2}{{A.4}{46}{Tema 4}{section*.52}{}}
\newlabel{sol:eje:4-inversa-z2@cref}{{[subappendix][4][2147483647,1]A.4}{[1][45][]46}}
\newlabel{sol:eje:4-comp-sobre}{{A.4}{46}{Tema 4}{section*.53}{}}
\newlabel{sol:eje:4-comp-sobre@cref}{{[subappendix][4][2147483647,1]A.4}{[1][46][]46}}
\newlabel{sol:eje:4-cambridge}{{A.4}{46}{Tema 4}{section*.54}{}}
\newlabel{sol:eje:4-cambridge@cref}{{[subappendix][4][2147483647,1]A.4}{[1][46][]46}}
\newlabel{sol:eje:4-iteradas}{{A.4}{47}{Tema 4}{section*.55}{}}
\newlabel{sol:eje:4-iteradas@cref}{{[subappendix][4][2147483647,1]A.4}{[1][47][]47}}
\newlabel{sol:eje:4-inv-comp}{{A.4}{47}{Tema 4}{section*.56}{}}
\newlabel{sol:eje:4-inv-comp@cref}{{[subappendix][4][2147483647,1]A.4}{[1][47][]47}}
\newlabel{sol:eje:4-derecha}{{A.4}{47}{Tema 4}{section*.57}{}}
\newlabel{sol:eje:4-derecha@cref}{{[subappendix][4][2147483647,1]A.4}{[1][47][]47}}
\newlabel{sol:eje:4-involucion}{{A.4}{48}{Tema 4}{section*.58}{}}
\newlabel{sol:eje:4-involucion@cref}{{[subappendix][4][2147483647,1]A.4}{[1][48][]48}}
\newlabel{sol:eje:4-cancelar}{{A.4}{48}{Tema 4}{section*.59}{}}
\newlabel{sol:eje:4-cancelar@cref}{{[subappendix][4][2147483647,1]A.4}{[1][48][]48}}
\newlabel{sol:eje:4-capas}{{A.4}{48}{Tema 4}{section*.60}{}}
\newlabel{sol:eje:4-capas@cref}{{[subappendix][4][2147483647,1]A.4}{[1][48][]48}}
\newlabel{sol:eje:4-tokenizador}{{A.4}{49}{Tema 4}{section*.61}{}}
\newlabel{sol:eje:4-tokenizador@cref}{{[subappendix][4][2147483647,1]A.4}{[1][49][]49}}
\newlabel{sol:eje:5-propiedades}{{A.5}{49}{Tema 5}{section*.62}{}}
\newlabel{sol:eje:5-propiedades@cref}{{[subappendix][5][2147483647,1]A.5}{[1][49][]49}}
\newlabel{sol:eje:5-enteros}{{A.5}{50}{Tema 5}{section*.63}{}}
\newlabel{sol:eje:5-enteros@cref}{{[subappendix][5][2147483647,1]A.5}{[1][49][]50}}
\newlabel{sol:eje:5-contar}{{A.5}{50}{Tema 5}{section*.64}{}}
\newlabel{sol:eje:5-contar@cref}{{[subappendix][5][2147483647,1]A.5}{[1][50][]50}}
\newlabel{sol:eje:5-clases}{{A.5}{50}{Tema 5}{section*.65}{}}
\newlabel{sol:eje:5-clases@cref}{{[subappendix][5][2147483647,1]A.5}{[1][50][]50}}
\newlabel{sol:eje:5-particion}{{A.5}{51}{Tema 5}{section*.66}{}}
\newlabel{sol:eje:5-particion@cref}{{[subappendix][5][2147483647,1]A.5}{[1][51][]51}}
\newlabel{sol:eje:5-bell}{{A.5}{51}{Tema 5}{section*.67}{}}
\newlabel{sol:eje:5-bell@cref}{{[subappendix][5][2147483647,1]A.5}{[1][51][]51}}
\newlabel{sol:eje:5-modn}{{A.5}{51}{Tema 5}{section*.68}{}}
\newlabel{sol:eje:5-modn@cref}{{[subappendix][5][2147483647,1]A.5}{[1][51][]51}}
\newlabel{sol:eje:5-nucleo}{{A.5}{52}{Tema 5}{section*.69}{}}
\newlabel{sol:eje:5-nucleo@cref}{{[subappendix][5][2147483647,1]A.5}{[1][52][]52}}
\newlabel{sol:eje:5-divide}{{A.5}{52}{Tema 5}{section*.70}{}}
\newlabel{sol:eje:5-divide@cref}{{[subappendix][5][2147483647,1]A.5}{[1][52][]52}}
\newlabel{sol:eje:5-mas3}{{A.5}{52}{Tema 5}{section*.71}{}}
\newlabel{sol:eje:5-mas3@cref}{{[subappendix][5][2147483647,1]A.5}{[1][52][]52}}
\newlabel{sol:eje:5-fracciones}{{A.5}{53}{Tema 5}{section*.72}{}}
\newlabel{sol:eje:5-fracciones@cref}{{[subappendix][5][2147483647,1]A.5}{[1][53][]53}}
\newlabel{sol:eje:5-orden-partes}{{A.5}{53}{Tema 5}{section*.73}{}}
\newlabel{sol:eje:5-orden-partes@cref}{{[subappendix][5][2147483647,1]A.5}{[1][53][]53}}
\newlabel{sol:eje:5-hasse}{{A.5}{53}{Tema 5}{section*.74}{}}
\newlabel{sol:eje:5-hasse@cref}{{[subappendix][5][2147483647,1]A.5}{[1][53][]53}}
\newlabel{sol:eje:5-lex}{{A.5}{54}{Tema 5}{section*.75}{}}
\newlabel{sol:eje:5-lex@cref}{{[subappendix][5][2147483647,1]A.5}{[1][54][]54}}
\newlabel{sol:eje:5-error}{{A.5}{54}{Tema 5}{section*.76}{}}
\newlabel{sol:eje:5-error@cref}{{[subappendix][5][2147483647,1]A.5}{[1][54][]54}}
\newlabel{sol:eje:6-vf}{{A.6}{55}{Tema 6}{section*.77}{}}
\newlabel{sol:eje:6-vf@cref}{{[subappendix][6][2147483647,1]A.6}{[1][55][]55}}
\newlabel{sol:eje:6-textos}{{A.6}{55}{Tema 6}{section*.78}{}}
\newlabel{sol:eje:6-textos@cref}{{[subappendix][6][2147483647,1]A.6}{[1][55][]55}}
\newlabel{sol:eje:6-meses}{{A.6}{55}{Tema 6}{section*.79}{}}
\newlabel{sol:eje:6-meses@cref}{{[subappendix][6][2147483647,1]A.6}{[1][55][]55}}
\newlabel{sol:eje:6-bits}{{A.6}{56}{Tema 6}{section*.80}{}}
\newlabel{sol:eje:6-bits@cref}{{[subappendix][6][2147483647,1]A.6}{[1][56][]56}}
\newlabel{sol:eje:6-restos}{{A.6}{56}{Tema 6}{section*.81}{}}
\newlabel{sol:eje:6-restos@cref}{{[subappendix][6][2147483647,1]A.6}{[1][56][]56}}
\newlabel{sol:eje:6-puntos}{{A.6}{56}{Tema 6}{section*.82}{}}
\newlabel{sol:eje:6-puntos@cref}{{[subappendix][6][2147483647,1]A.6}{[1][56][]56}}
\newlabel{sol:eje:6-compresion}{{A.6}{57}{Tema 6}{section*.83}{}}
\newlabel{sol:eje:6-compresion@cref}{{[subappendix][6][2147483647,1]A.6}{[1][57][]57}}
\newlabel{sol:eje:6-pi}{{A.6}{57}{Tema 6}{section*.84}{}}
\newlabel{sol:eje:6-pi@cref}{{[subappendix][6][2147483647,1]A.6}{[1][57][]57}}
\newlabel{sol:eje:6-shortlex}{{A.6}{57}{Tema 6}{section*.85}{}}
\newlabel{sol:eje:6-shortlex@cref}{{[subappendix][6][2147483647,1]A.6}{[1][57][]57}}
\newlabel{sol:eje:6-diagonal}{{A.6}{58}{Tema 6}{section*.86}{}}
\newlabel{sol:eje:6-diagonal@cref}{{[subappendix][6][2147483647,1]A.6}{[1][57][]58}}
\newlabel{sol:eje:6-subconjN}{{A.6}{58}{Tema 6}{section*.87}{}}
\newlabel{sol:eje:6-subconjN@cref}{{[subappendix][6][2147483647,1]A.6}{[1][58][]58}}
\newlabel{sol:eje:6-producto}{{A.6}{58}{Tema 6}{section*.88}{}}
\newlabel{sol:eje:6-producto@cref}{{[subappendix][6][2147483647,1]A.6}{[1][58][]58}}
\newlabel{sol:eje:6-Q}{{A.6}{59}{Tema 6}{section*.89}{}}
\newlabel{sol:eje:6-Q@cref}{{[subappendix][6][2147483647,1]A.6}{[1][59][]59}}
\newlabel{sol:eje:6-PA}{{A.6}{59}{Tema 6}{section*.90}{}}
\newlabel{sol:eje:6-PA@cref}{{[subappendix][6][2147483647,1]A.6}{[1][59][]59}}
\newlabel{sol:eje:6-intervalos}{{A.6}{59}{Tema 6}{section*.91}{}}
\newlabel{sol:eje:6-intervalos@cref}{{[subappendix][6][2147483647,1]A.6}{[1][59][]59}}
\makeatother

\begin{document}
\frontmatter
\begin{titlepage}
\centering
\vspace*{2.5cm}
{\large\scshape Anatomía de un LLM}\par
\vspace{0.4cm}
{\normalsize Fundamentos matemáticos de los modelos de lenguaje}\par
\vspace{2.2cm}
\rule{\textwidth}{0.8pt}\par\vspace{0.5cm}
{\small\scshape Fase 0 \textperiodcentered\ Preliminares matemáticos}\par\vspace{0.3cm}
{\Huge\bfseries Módulo 0.2\par}\vspace{0.3cm}
{\LARGE Conjuntos, funciones y relaciones\par}
\vspace{0.5cm}\rule{\textwidth}{0.4pt}\par
\vspace{1.6cm}
{\large Resumen del módulo}\par
\vfill
{\small Edición de octubre de 2026}\par
\end{titlepage}

\chapter*{Sobre este resumen}
\addcontentsline{toc}{chapter}{Sobre este resumen}
\markboth{Sobre este resumen}{}

Este documento reúne lo indispensable de los apuntes del Módulo 0.2: el texto expositivo, todas las definiciones y todos los resultados (teoremas, proposiciones, lemas y corolarios), las tablas y las observaciones de contenido. Se han suprimido las demostraciones, los ejemplos, las lecturas, las figuras, los ejercicios y sus soluciones. La numeración de temas, secciones, enunciados y ecuaciones es la misma que en los apuntes, de modo que cualquier referencia puede buscarse allí; las que remiten a elementos suprimidos (un ejemplo, una figura, una ecuación de una demostración) conservan su número en los apuntes.

\section*{Notación}
\begin{center}
\begin{tabular}{@{}ll@{}}
\toprule
Símbolo & Significado \\
\midrule
$\N=\{0,1,2,\dots\}$ & números naturales (incluido el~0) \\
$\Z$, $\Q$, $\R$ & enteros, racionales, reales \\
$\abs{A}$ & número de elementos (cardinal) de un conjunto finito $A$ \\
$\Pa(A)$ & conjunto de las partes de $A$ \\
$A\setminus B$, $\overline{A}$, $A\dsim B$ & diferencia, complemento, diferencia simétrica \\
$f\colon A\to B$ & función de dominio $A$ y codominio $B$ \\
$f(X)$, $f^{-1}(Y)$ & imagen de $X\subseteq A$ y preimagen de $Y\subseteq B$ \\
$g\circ f$ & composición: primero $f$, después $g$ \\
$\id_A$ & función identidad de $A$ \\
$B^{A}$ & conjunto de todas las funciones $A\to B$ \\
$\Sigma^{*}$, $\vacia$, $\abs{w}$ & cadenas sobre el alfabeto $\Sigma$, cadena vacía, longitud \\
$[a]$, $A/{\sim}$ & clase de equivalencia y conjunto cociente \\
\bottomrule
\end{tabular}
\end{center}


\tableofcontents

\mainmatter
% =====================================================================
\chapter{Conjuntos y operaciones}
\label{tema:1}
% =====================================================================

Este tema fija el vocabulario básico de la teoría de conjuntos y, sobre todo, la técnica para demostrar que dos conjuntos son iguales. La presentación sigue el capítulo~1 y el capítulo~8 de Hammack~\cite{hammack}, las secciones~1.3, 1.4 y~2.3 de Velleman~\cite{velleman} y la sección~4.1 de~\cite{mcs}; el estilo de las demostraciones es el de la clase «Set Theory Revisited» de Stanford CS103~\cite{cs103set}.

\section{Conjuntos y pertenencia}

\begin{definicion}[Conjunto]
Un \emph{conjunto} es una colección de objetos, llamados \emph{elementos}, considerada como un objeto en sí misma. Se escribe $x\in A$ si $x$ es elemento de $A$ y $x\notin A$ en caso contrario. Dos conjuntos son \emph{iguales} si tienen exactamente los mismos elementos:
\[
A=B\iff \forall x\,(x\in A\leftrightarrow x\in B).
\]
\end{definicion}

Un conjunto puede describirse \emph{por extensión}, enumerando sus elementos, como $\{1,2,3\}$, o \emph{por comprensión}, mediante una propiedad, como $\{x\in\Z : x^{2}<5\}=\{-2,-1,0,1,2\}$. Por la definición de igualdad, el orden y las repeticiones no importan: $\{1,2,2\}=\{2,1\}$. El \emph{conjunto vacío} $\emptyset$ no tiene elementos. Si $A$ tiene un número finito de elementos, ese número es su \emph{cardinal}, $\abs{A}$.

\addtocounter{definicion}{1}%

\section{Subconjuntos y conjunto de partes}

\begin{definicion}[Subconjunto]
$A$ es \emph{subconjunto} de $B$, y se escribe $A\subseteq B$, si todo elemento de $A$ es elemento de $B$: $\forall x\,(x\in A\rightarrow x\in B)$. Si además $A\neq B$, se dice que $A$ es un subconjunto \emph{propio} y se escribe $A\subsetneq B$.
\end{definicion}

\begin{proposicion}\label{prop:t1-subconj}
Para conjuntos cualesquiera $A$, $B$ y $C$:
\begin{enumerate}
  \item $\emptyset\subseteq A$ y $A\subseteq A$;
  \item si $A\subseteq B$ y $B\subseteq C$, entonces $A\subseteq C$;
  \item $A=B$ si y solo si $A\subseteq B$ y $B\subseteq A$.
\end{enumerate}
\end{proposicion}

El apartado~(iii) es el \emph{método de la doble inclusión}: para probar $A=B$ se toma un elemento arbitrario de $A$ y se demuestra que está en $B$, y después al revés~\cite[§8.3]{hammack}.

\begin{definicion}[Conjunto de partes]
El \emph{conjunto de las partes} de $A$ es el conjunto de todos sus subconjuntos: $\Pa(A)=\{X : X\subseteq A\}$.
\end{definicion}

\begin{teorema}\label{teo:t1-partes}
Si $A$ es finito con $\abs{A}=n$, entonces $\abs{\Pa(A)}=2^{n}$.
\end{teorema}

\addtocounter{definicion}{1}%

\section{Operaciones con conjuntos}

\begin{definicion}[Operaciones]
Sean $A$ y $B$ subconjuntos de un conjunto universal $U$. Se definen
\begin{align*}
A\cup B&=\{x : x\in A\ \text{o}\ x\in B\} &&\text{(unión)},\\
A\cap B&=\{x : x\in A\ \text{y}\ x\in B\} &&\text{(intersección)},\\
A\setminus B&=\{x : x\in A\ \text{y}\ x\notin B\} &&\text{(diferencia)},\\
\overline{A}&=U\setminus A &&\text{(complemento)},\\
A\dsim B&=(A\setminus B)\cup(B\setminus A) &&\text{(diferencia simétrica)}.
\end{align*}
$A$ y $B$ son \emph{disjuntos} si $A\cap B=\emptyset$.
\end{definicion}

Cada operación corresponde a un conector lógico: la unión a la disyunción, la intersección a la conjunción y el complemento a la negación. Por eso las leyes de la lógica proposicional se traducen en leyes de conjuntos~\cite[§1.4]{velleman}.

\begin{teorema}[Leyes del álgebra de conjuntos]\label{teo:t1-leyes}
Para subconjuntos $A$, $B$, $C$ de $U$:
\begin{enumerate}
  \item conmutativas: $A\cup B=B\cup A$, \ $A\cap B=B\cap A$;
  \item asociativas: $(A\cup B)\cup C=A\cup(B\cup C)$, \ $(A\cap B)\cap C=A\cap(B\cap C)$;
  \item distributivas: $A\cap(B\cup C)=(A\cap B)\cup(A\cap C)$, \ $A\cup(B\cap C)=(A\cup B)\cap(A\cup C)$;
  \item leyes de De Morgan: $\overline{A\cup B}=\overline{A}\cap\overline{B}$, \ $\overline{A\cap B}=\overline{A}\cup\overline{B}$;
  \item $A\setminus B=A\cap\overline{B}$ \ y \ $\overline{\overline{A}}=A$.
\end{enumerate}
\end{teorema}

\addtocounter{figure}{1}%

\addtocounter{definicion}{1}%

\addtocounter{definicion}{1}%

\section{Producto cartesiano y familias indexadas}

\begin{definicion}[Producto cartesiano]
Un \emph{par ordenado} $(a,b)$ es un objeto con primera componente $a$ y segunda $b$, con la propiedad $(a,b)=(c,d)\iff a=c\ \text{y}\ b=d$. El \emph{producto cartesiano} de $A$ y $B$ es $A\times B=\{(a,b): a\in A,\ b\in B\}$. Análogamente, $A_1\times\dots\times A_n$ es el conjunto de $n$-tuplas, y $A^{n}=A\times\dots\times A$ ($n$ factores).
\end{definicion}

\begin{proposicion}\label{prop:t1-producto}
Si $A$ y $B$ son finitos, $\abs{A\times B}=\abs{A}\cdot\abs{B}$; en particular, $\abs{A^{n}}=\abs{A}^{n}$. Además, para conjuntos cualesquiera, $A\times(B\cup C)=(A\times B)\cup(A\times C)$.
\end{proposicion}

\addtocounter{definicion}{1}%

\begin{definicion}[Familias indexadas]
Sea $I$ un conjunto de índices y, para cada $i\in I$, un conjunto $A_i$. Se definen
\[
\bigcup_{i\in I}A_i=\{x : x\in A_i\ \text{para algún}\ i\in I\},\qquad
\bigcap_{i\in I}A_i=\{x : x\in A_i\ \text{para todo}\ i\in I\}.
\]
\end{definicion}

\addtocounter{definicion}{1}%

\addtocounter{definicion}{1}%


% =====================================================================
\chapter{Funciones, imagen y preimagen}
\label{tema:2}
% =====================================================================

Este tema define con precisión qué es una función y estudia cómo actúa sobre subconjuntos mediante la imagen y la preimagen. Sigue las secciones~12.1 y~12.6 de Hammack~\cite{hammack}, las secciones~5.1 y~5.5 de Velleman~\cite{velleman} y la sección~4.3 de~\cite{mcs}. Las propiedades de la imagen y la preimagen son las de la pregunta~9 de la hoja~1 de Cambridge~\cite{camns}.

\section{Definición de función}

\begin{definicion}[Función]\label{def:t2-funcion}
Sean $A$ y $B$ conjuntos. Una \emph{función} $f$ de $A$ en $B$, escrita $f\colon A\to B$, es un subconjunto $f\subseteq A\times B$ tal que para cada $a\in A$ existe \emph{exactamente un} $b\in B$ con $(a,b)\in f$; ese $b$ se denota $f(a)$. El conjunto $A$ es el \emph{dominio} y $B$ el \emph{codominio}. El \emph{recorrido} (o \emph{imagen}) de $f$ es $f(A)=\{f(a): a\in A\}\subseteq B$.
\end{definicion}

La definición contiene dos exigencias: \emph{existencia} (todo elemento del dominio tiene imagen) y \emph{unicidad} (no tiene dos). Dos funciones $f,g\colon A\to B$ son iguales si $f(a)=g(a)$ para todo $a\in A$; la igualdad exige además que el dominio y el codominio coincidan.

\addtocounter{definicion}{1}%

\addtocounter{definicion}{1}%

\section{El conjunto de funciones}

\begin{definicion}
$B^{A}$ denota el conjunto de todas las funciones $A\to B$.
\end{definicion}

\begin{proposicion}\label{prop:t2-contar}
Si $A$ y $B$ son finitos, $\abs{B^{A}}=\abs{B}^{\abs{A}}$.
\end{proposicion}

\begin{definicion}[Función característica]
Sea $X\subseteq U$. La \emph{función característica} (o \emph{indicadora}) de $X$ es $\chi_X\colon U\to\{0,1\}$, con $\chi_X(u)=1$ si $u\in X$ y $\chi_X(u)=0$ si $u\notin X$.
\end{definicion}

Cada subconjunto de $U$ determina su función característica y cada función $U\to\{0,1\}$ es la característica de un único subconjunto, el de los elementos que van al~1. Con $\abs{U}=n$, la \cref{prop:t2-contar} da $\abs{\{0,1\}^{U}}=2^{n}$, de nuevo el \cref{teo:t1-partes}. En el \cref{tema:3} se formaliza esta correspondencia como una biyección.

\section{Imagen y preimagen de subconjuntos}

\begin{definicion}[Imagen y preimagen]\label{def:t2-imagen}
Sea $f\colon A\to B$. Para $X\subseteq A$ y $Y\subseteq B$ se definen la \emph{imagen} de $X$ y la \emph{preimagen} (o \emph{imagen inversa}) de $Y$:
\[
f(X)=\{f(x): x\in X\}\subseteq B,\qquad f^{-1}(Y)=\{a\in A: f(a)\in Y\}\subseteq A.
\]
\end{definicion}

La notación $f^{-1}(Y)$ tiene sentido para \emph{cualquier} función, sea o no invertible: es un conjunto, no la aplicación de una función inversa. Si $Y=\{b\}$, el conjunto $f^{-1}(\{b\})$ se llama \emph{fibra} de $b$.

\addtocounter{definicion}{1}%

\begin{teorema}[Propiedades de la preimagen]\label{teo:t2-preimagen}
Sea $f\colon A\to B$ y sean $C,D\subseteq B$. Entonces
\begin{gather*}
f^{-1}(C\cup D)=f^{-1}(C)\cup f^{-1}(D),\qquad f^{-1}(C\cap D)=f^{-1}(C)\cap f^{-1}(D),\\
f^{-1}(B\setminus C)=A\setminus f^{-1}(C).
\end{gather*}
\end{teorema}

\begin{teorema}[Propiedades de la imagen]\label{teo:t2-imagen}
Sea $f\colon A\to B$, $X_1,X_2,X\subseteq A$ e $Y\subseteq B$. Entonces
\begin{enumerate}
  \item $f(X_1\cup X_2)=f(X_1)\cup f(X_2)$;
  \item $f(X_1\cap X_2)\subseteq f(X_1)\cap f(X_2)$, y la inclusión puede ser estricta;
  \item $X\subseteq f^{-1}(f(X))$ y $f(f^{-1}(Y))\subseteq Y$.
\end{enumerate}
\end{teorema}

La asimetría entre ambos teoremas es instructiva: la preimagen respeta todas las operaciones de conjuntos, mientras que la imagen solo respeta la unión. La razón es que la imagen puede «pegar» elementos distintos; en el \cref{tema:3} se verá que las igualdades $f(X_1\cap X_2)=f(X_1)\cap f(X_2)$ y $f^{-1}(f(X))=X$ son ciertas exactamente cuando $f$ es inyectiva.

\section{Otras construcciones}

Una \emph{sucesión} de elementos de $A$ es una función $\N\to A$, $n\mapsto a_n$. Una \emph{cadena} de longitud $n$ sobre $\Sigma$ puede verse como una función $\{1,\dots,n\}\to\Sigma$ que asigna a cada posición su símbolo. Si $f\colon A\to B$ y $X\subseteq A$, la \emph{restricción} $f|_X\colon X\to B$ es la función $f|_X(x)=f(x)$. Las funciones \emph{definidas a trozos} asignan reglas distintas a partes disjuntas del dominio que lo recubren.

\addtocounter{definicion}{1}%

\addtocounter{definicion}{1}%


% =====================================================================
\chapter{Funciones inyectivas, sobreyectivas y biyectivas}
\label{tema:3}
% =====================================================================

Este tema es el núcleo práctico del módulo: decidir si una función es inyectiva, sobreyectiva o biyectiva y \emph{demostrarlo}. Sigue la sección~12.2 de Hammack~\cite{hammack}, la sección~5.2 de Velleman~\cite{velleman} y las secciones~4.3 y~4.5 de~\cite{mcs}. Los ejercicios de clasificación siguen el problema~1 de la hoja~3 de Stanford CS103~\cite{cs103ps3}.

\section{Definiciones y método de demostración}

\begin{definicion}\label{def:t3-iny}
Sea $f\colon A\to B$.
\begin{enumerate}
  \item $f$ es \emph{inyectiva} si elementos distintos tienen imágenes distintas: para todos $a,a'\in A$, si $f(a)=f(a')$ entonces $a=a'$.
  \item $f$ es \emph{sobreyectiva} si todo elemento del codominio es imagen de alguno: para todo $b\in B$ existe $a\in A$ con $f(a)=b$; equivalentemente, $f(A)=B$.
  \item $f$ es \emph{biyectiva} si es inyectiva y sobreyectiva.
\end{enumerate}
\end{definicion}

En términos de fibras: $f$ es inyectiva si cada fibra $f^{-1}(\{b\})$ tiene \emph{a lo sumo} un elemento, sobreyectiva si tiene \emph{al menos} uno y biyectiva si tiene \emph{exactamente} uno. La \cref{fig:t3-flechas} muestra los cuatro casos posibles con diagramas de flechas.

\addtocounter{figure}{1}%

\paragraph{Cómo demostrar cada propiedad.}
\begin{itemize}[leftmargin=1.4em]
  \item \emph{Inyectividad}: se supone $f(a)=f(a')$ para $a,a'$ arbitrarios y se deduce $a=a'$ (demostración directa), o se supone $a\neq a'$ y se deduce $f(a)\neq f(a')$ (contrarrecíproco).
  \item \emph{No inyectividad}: basta un contraejemplo, $a\neq a'$ con $f(a)=f(a')$.
  \item \emph{Sobreyectividad}: se toma $b\in B$ arbitrario, se \emph{busca} un candidato $a$ (normalmente resolviendo la ecuación $f(a)=b$ en borrador), y se \emph{comprueba} que $a\in A$ y que $f(a)=b$.
  \item \emph{No sobreyectividad}: basta un $b\in B$ y una prueba de que la ecuación $f(a)=b$ no tiene solución en $A$.
\end{itemize}

\addtocounter{definicion}{1}%

\addtocounter{definicion}{1}%

\addtocounter{definicion}{1}%

\addtocounter{definicion}{1}%

\section{Criterios para funciones reales}

\begin{definicion}
$f\colon A\to\R$, con $A\subseteq\R$, es \emph{estrictamente creciente} si $x<x'$ implica $f(x)<f(x')$, y \emph{estrictamente decreciente} si $x<x'$ implica $f(x)>f(x')$. Es \emph{estrictamente monótona} si es una de las dos.
\end{definicion}

\begin{proposicion}\label{prop:t3-monotona}
Toda función estrictamente monótona es inyectiva.
\end{proposicion}

Si $f$ es derivable en un intervalo con $f'>0$ en todo punto, entonces es estrictamente creciente (consecuencia del teorema del valor medio, que se estudia en el módulo de cálculo de esta fase), y por tanto inyectiva. Para la sobreyectividad sobre un intervalo, el teorema de Bolzano (valores intermedios) permite concluir que una función continua toma todos los valores comprendidos entre sus límites; en este módulo se prefiere resolver $f(x)=y$ explícitamente.

\addtocounter{definicion}{1}%

\section{Funciones entre conjuntos finitos}

\begin{proposicion}\label{prop:t3-finitos}
Sean $A$ y $B$ finitos y $f\colon A\to B$.
\begin{enumerate}
  \item Si $f$ es inyectiva, $\abs{A}\le\abs{B}$; si es sobreyectiva, $\abs{A}\ge\abs{B}$.
  \item Si $\abs{A}=\abs{B}$, entonces $f$ es inyectiva si y solo si es sobreyectiva.
\end{enumerate}
\end{proposicion}

\begin{proposicion}\label{prop:t3-contar-iny}
Si $\abs{A}=k$ y $\abs{B}=n$, el número de funciones inyectivas $A\to B$ es $n(n-1)\cdots(n-k+1)$, que vale~0 si $k>n$. En particular, hay $n!$ biyecciones entre dos conjuntos de $n$ elementos.
\end{proposicion}

\addtocounter{definicion}{1}%

\addtocounter{definicion}{1}%


% =====================================================================
\chapter{Composición e inversa}
\label{tema:4}
% =====================================================================

Componer funciones es aplicarlas una tras otra; invertir una función es deshacerla. Este tema estudia ambas operaciones y cómo se comportan respecto de la inyectividad y la sobreyectividad. Sigue las secciones~12.4 y~12.5 de Hammack~\cite{hammack}, las secciones~5.2 y~5.3 de Velleman~\cite{velleman} y la sección~4.3 de~\cite{mcs}.

\section{Composición}

\begin{definicion}[Composición]
Sean $f\colon A\to B$ y $g\colon B\to C$. La \emph{composición} de $g$ con $f$ es la función $g\circ f\colon A\to C$ definida por $(g\circ f)(a)=g(f(a))$.
\end{definicion}

El orden de lectura es de derecha a izquierda: en $g\circ f$ se aplica primero $f$. La composición solo está definida cuando el codominio de $f$ está contenido en el dominio de $g$.

\addtocounter{definicion}{1}%

\begin{proposicion}\label{prop:t4-asoc}
Sean $f\colon A\to B$, $g\colon B\to C$ y $h\colon C\to D$. Entonces $h\circ(g\circ f)=(h\circ g)\circ f$, y $f\circ\id_A=f=\id_B\circ f$.
\end{proposicion}

\section{Composición, inyectividad y sobreyectividad}

\begin{teorema}\label{teo:t4-comp}
Sean $f\colon A\to B$ y $g\colon B\to C$.
\begin{enumerate}
  \item Si $f$ y $g$ son inyectivas, $g\circ f$ es inyectiva.
  \item Si $f$ y $g$ son sobreyectivas, $g\circ f$ es sobreyectiva.
  \item Si $f$ y $g$ son biyectivas, $g\circ f$ es biyectiva.
\end{enumerate}
\end{teorema}

\begin{proposicion}\label{prop:t4-reciproco}
Sean $f\colon A\to B$ y $g\colon B\to C$.
\begin{enumerate}
  \item Si $g\circ f$ es inyectiva, $f$ es inyectiva.
  \item Si $g\circ f$ es sobreyectiva, $g$ es sobreyectiva.
\end{enumerate}
\end{proposicion}

Las otras dos implicaciones son falsas: si $g\circ f$ es inyectiva, $g$ no tiene por qué serlo, y si $g\circ f$ es sobreyectiva, $f$ no tiene por qué serlo (\cref{eje:4-cambridge}).

\section{Función inversa}

\begin{definicion}[Inversa]
Sea $f\colon A\to B$. Una función $g\colon B\to A$ es \emph{inversa} de $f$ si $g\circ f=\id_A$ y $f\circ g=\id_B$. Si solo se cumple $g\circ f=\id_A$, se dice que $g$ es una \emph{inversa por la izquierda} de $f$; si solo se cumple $f\circ g=\id_B$, una \emph{inversa por la derecha}.
\end{definicion}

\begin{teorema}\label{teo:t4-inversa}
Una función $f\colon A\to B$ tiene inversa si y solo si es biyectiva. En ese caso la inversa es única; se denota $f^{-1}$ y está dada por $f^{-1}(b)=$ el único $a\in A$ con $f(a)=b$.
\end{teorema}

\begin{proposicion}\label{prop:t4-laterales}
Sea $f\colon A\to B$ con $A\neq\emptyset$. Entonces $f$ tiene una inversa por la izquierda si y solo si es inyectiva, y tiene una inversa por la derecha si y solo si es sobreyectiva.
\end{proposicion}

\begin{proposicion}\label{prop:t4-inv-comp}
Si $f\colon A\to B$ es biyectiva, $f^{-1}$ es biyectiva y $(f^{-1})^{-1}=f$. Si además $g\colon B\to C$ es biyectiva, $(g\circ f)^{-1}=f^{-1}\circ g^{-1}$.
\end{proposicion}

\section{Cálculo de inversas}

Para invertir una biyección real $y=f(x)$ se despeja $x$ en función de $y$ y se comprueba que el resultado pertenece al dominio. Cuando $f$ no es biyectiva, puede \emph{restringirse} el dominio y el codominio para obtener una biyección: así se definen $\sqrt{\phantom{x}}$ como inversa de $x^{2}$ en $[0,\infty)$, $\ln$ como inversa de $\exp\colon\R\to(0,\infty)$ y $\operatorname{arcsen}$ como inversa de $\operatorname{sen}$ restringido a $[-\pi/2,\pi/2]$.

\addtocounter{definicion}{1}%

\addtocounter{definicion}{1}%

\addtocounter{definicion}{1}%


% =====================================================================
\chapter{Relaciones de equivalencia y de orden}
\label{tema:5}
% =====================================================================

Una relación binaria expresa que dos elementos están «relacionados» de cierta manera: ser iguales módulo~$n$, tener la misma longitud, preceder a otro. Este tema estudia las dos familias de relaciones más importantes, las de equivalencia, que clasifican, y las de orden, que ordenan. Sigue el capítulo~11 de Hammack~\cite{hammack}, las secciones~4.2 a~4.5 de Velleman~\cite{velleman} y las secciones~4.4 y~9.6 a~9.10 de~\cite{mcs}.

\section{Relaciones binarias}

\begin{definicion}[Relación]
Una \emph{relación} (binaria) en un conjunto $A$ es un subconjunto $R\subseteq A\times A$. Se escribe $a\,R\,b$ si $(a,b)\in R$.
\end{definicion}

\addtocounter{definicion}{1}%

Una relación en un conjunto finito puede representarse con un grafo dirigido, con un arco $a\to b$ por cada par $(a,b)\in R$, o con una matriz de ceros y unos cuya entrada $(a,b)$ vale~1 si $a\,R\,b$.

\begin{definicion}[Propiedades]\label{def:t5-prop}
Una relación $R$ en $A$ es
\begin{enumerate}
  \item \emph{reflexiva} si $a\,R\,a$ para todo $a\in A$;
  \item \emph{simétrica} si $a\,R\,b$ implica $b\,R\,a$;
  \item \emph{antisimétrica} si $a\,R\,b$ y $b\,R\,a$ implican $a=b$;
  \item \emph{transitiva} si $a\,R\,b$ y $b\,R\,c$ implican $a\,R\,c$.
\end{enumerate}
\end{definicion}

\begin{center}\small
\begin{tabular}{@{}lcccc@{}}
\toprule
Relación & Reflexiva & Simétrica & Antisimétrica & Transitiva \\
\midrule
$=$ en $\Z$ & sí & sí & sí & sí \\
$\le$ en $\Z$ & sí & no & sí & sí \\
$<$ en $\Z$ & no & no & sí & sí \\
$\mid$ en $\N$ & sí & no & sí & sí \\
$\equiv$ (mód.\ $n$) en $\Z$ & sí & sí & no & sí \\
$\neq$ en $\Z$ & no & sí & no & no \\
\bottomrule
\end{tabular}
\end{center}

Simetría y antisimetría \emph{no} son contrarias: la igualdad tiene ambas, y la relación $\{(1,2),(2,1),(1,3)\}$ en $\{1,2,3\}$ no tiene ninguna.

\section{Relaciones de equivalencia y particiones}

\begin{definicion}[Relación de equivalencia]
Una relación $\sim$ en $A$ es de \emph{equivalencia} si es reflexiva, simétrica y transitiva. La \emph{clase de equivalencia} de $a\in A$ es $[a]=\{x\in A: x\sim a\}$, y el \emph{conjunto cociente} es $A/{\sim}=\{[a]: a\in A\}$.
\end{definicion}

\begin{definicion}[Partición]
Una \emph{partición} de $A$ es una familia de subconjuntos de $A$ no vacíos, disjuntos dos a dos y cuya unión es $A$.
\end{definicion}

\begin{teorema}\label{teo:t5-particion}
Sea $\sim$ una relación de equivalencia en $A$. Entonces:
\begin{enumerate}
  \item $a\in[a]$ para todo $a$;
  \item $[a]=[b]$ si y solo si $a\sim b$;
  \item dos clases son iguales o disjuntas.
\end{enumerate}
En consecuencia, $A/{\sim}$ es una partición de $A$. Recíprocamente, toda partición $\mathcal{C}$ de $A$ define una relación de equivalencia, $a\sim b$ si $a$ y $b$ están en el mismo bloque de $\mathcal{C}$, cuyas clases son los bloques.
\end{teorema}

\addtocounter{definicion}{1}%

\begin{proposicion}[Núcleo de una función]\label{prop:t5-nucleo}
Sea $f\colon A\to B$. La relación $a\sim_f a'$ si $f(a)=f(a')$ es de equivalencia, y sus clases son las fibras no vacías $f^{-1}(\{b\})$, $b\in f(A)$.
\end{proposicion}

La demostración es el \cref{eje:5-nucleo}. Muchas relaciones de equivalencia son de este tipo: «tener la misma longitud» es el núcleo de $\ell\colon\Sigma^{*}\to\N$, y la congruencia módulo $n$ es el núcleo de la función «resto de dividir entre $n$».

\section{Relaciones de orden}

\begin{definicion}[Orden]
Una relación $\preceq$ en $A$ es un \emph{orden parcial} si es reflexiva, antisimétrica y transitiva; el par $(A,\preceq)$ es un \emph{conjunto ordenado}. Es un \emph{orden total} si, además, dos elementos cualesquiera son comparables: $a\preceq b$ o $b\preceq a$. Un elemento $m$ es \emph{minimal} si no existe $x\neq m$ con $x\preceq m$, y \emph{mínimo} si $m\preceq x$ para todo $x$; análogamente se definen maximal y máximo.
\end{definicion}

\addtocounter{definicion}{1}%

Un conjunto ordenado finito se representa con su \emph{diagrama de Hasse}: se dibuja un segmento ascendente de $x$ a $y$ cuando $x\prec y$ y no hay ningún $z$ con $x\prec z\prec y$; la reflexividad y la transitividad quedan implícitas (\cref{fig:t5-hasse}).

\addtocounter{figure}{1}%

\addtocounter{definicion}{1}%

\addtocounter{definicion}{1}%


% =====================================================================
\chapter{Cardinalidad, palomar y diagonal de Cantor}
\label{tema:6}
% =====================================================================

Las funciones inyectivas, sobreyectivas y biyectivas permiten comparar el tamaño de dos conjuntos sin contar sus elementos. En conjuntos finitos esta comparación da el principio del palomar; en conjuntos infinitos conduce a la distinción entre conjuntos numerables y no numerables y al argumento diagonal de Cantor. El tema sigue la sección~12.3 y el capítulo~14 de Hammack~\cite{hammack}, el capítulo~8 de Velleman~\cite{velleman}, las secciones~4.5, 7.1 y~14.8 de~\cite{mcs} y las transparencias «Cardinality» de Stanford CS103~\cite{cs103card}.

\section{Comparar tamaños con funciones}

\begin{definicion}[Equipotencia]
Dos conjuntos $A$ y $B$ tienen el \emph{mismo cardinal}, y se escribe $\abs{A}=\abs{B}$, si existe una biyección $A\to B$. Se escribe $\abs{A}\le\abs{B}$ si existe una función inyectiva $A\to B$.
\end{definicion}

Para conjuntos finitos, estas definiciones coinciden con contar elementos (\cref{prop:t3-finitos}). La relación «tener el mismo cardinal» es reflexiva ($\id_A$ es biyectiva), simétrica (la inversa de una biyección es biyectiva, \cref{prop:t4-inv-comp}) y transitiva (la composición de biyecciones es biyectiva, \cref{teo:t4-comp}).

\addtocounter{definicion}{1}%

\section{El principio del palomar}

\begin{teorema}[Principio del palomar]\label{teo:t6-palomar}
Sean $A$ y $B$ finitos y $f\colon A\to B$.
\begin{enumerate}
  \item Si $\abs{A}>\abs{B}$, $f$ no es inyectiva: existen $a\neq a'$ con $f(a)=f(a')$.
  \item (Forma general.) Si $\abs{A}>k\abs{B}$ con $k\in\N$, alguna fibra tiene al menos $k+1$ elementos. En particular, alguna fibra tiene al menos $\lceil\abs{A}/\abs{B}\rceil$ elementos.
\end{enumerate}
\end{teorema}

El nombre procede de la imagen de $\abs{A}$ palomas que se refugian en $\abs{B}$ palomares: si hay más palomas que palomares, en algún palomar hay al menos dos. La dificultad en las aplicaciones está en elegir bien las «palomas» y los «palomares».

\addtocounter{definicion}{1}%

\begin{proposicion}[Límite de la compresión sin pérdida]\label{prop:t6-compresion}
Para $n\ge1$ no existe ninguna función inyectiva de $\{0,1\}^{n}$ en el conjunto $\{0,1\}^{<n}=\bigcup_{k=0}^{n-1}\{0,1\}^{k}$ de cadenas binarias de longitud menor que $n$.
\end{proposicion}

Un compresor sin pérdida es una función inyectiva (debe poder descomprimirse). La \cref{prop:t6-compresion} dice que ningún compresor acorta \emph{todas} las cadenas de longitud $n$: si acorta algunas, alarga o deja igual otras. El módulo~6 cuantifica esta idea con la entropía.

\section{Conjuntos numerables}

\begin{definicion}[Numerable]
Un conjunto $A$ es \emph{infinito numerable} si $\abs{A}=\abs{\N}$, es decir, si sus elementos pueden listarse como $a_0,a_1,a_2,\dots$ sin repeticiones y sin omitir ninguno. Es \emph{numerable} si es finito o infinito numerable, y \emph{no numerable} en otro caso.
\end{definicion}

\begin{proposicion}\label{prop:t6-numerables}
Los conjuntos $\Z$, $\N\times\N$ y $\Sigma^{*}$ (con $\Sigma$ finito y no vacío) son infinitos numerables.
\end{proposicion}

\addtocounter{figure}{1}%

\begin{proposicion}\label{prop:t6-criterios}
Sea $A$ un conjunto no vacío. Son equivalentes: (i) $A$ es numerable; (ii) existe una función inyectiva $A\to\N$; (iii) existe una función sobreyectiva $\N\to A$. Además, todo subconjunto de un conjunto numerable es numerable, y la unión de una familia numerable de conjuntos numerables es numerable.
\end{proposicion}

Se admite este resultado sin demostración completa (véase~\cite[§14.2]{hammack} y~\cite[§8.2]{velleman}); algunas de sus partes se piden en los \cref{eje:6-subconjN,eje:6-producto}. Como consecuencia, $\Q$ es numerable (\cref{eje:6-Q}).

\section{El argumento diagonal de Cantor}

\begin{teorema}[Cantor]\label{teo:t6-cantor}
Para todo conjunto $A$ no existe ninguna función sobreyectiva $A\to\Pa(A)$. En consecuencia, $\abs{A}\neq\abs{\Pa(A)}$.
\end{teorema}

\begin{corolario}\label{cor:t6-nonum}
$\Pa(\N)$ y el conjunto $\{0,1\}^{\N}$ de sucesiones infinitas de ceros y unos no son numerables.
\end{corolario}

El nombre «diagonal» se entiende con sucesiones: dada una lista $s_0,s_1,s_2,\dots$ de sucesiones binarias, la sucesión $d$ con $d_k=1-s_k(k)$ difiere de $s_k$ en la posición $k$, luego no está en la lista. El mismo argumento, con desarrollos decimales y evitando las colas de nueves, prueba que $\R$ no es numerable~\cite[§14.2]{hammack}.

\addtocounter{definicion}{1}%

\addtocounter{definicion}{1}%



\backmatter
% =====================================================================
\begin{thebibliography}{99}
\addcontentsline{toc}{chapter}{Bibliografía}
\markboth{Bibliografía}{}
\small

\bibitem{camdm}
University of Cambridge, Department of Computer Science and Technology.
\newblock Discrete Mathematics, curso 2021--22.
\newblock Hoja de ejercicios~8: «On surjections and injections», «On images», «On countability».
\newblock \url{https://www.cl.cam.ac.uk/teaching/2122/DiscMath/}

\bibitem{camns}
University of Cambridge, Faculty of Mathematics.
\newblock IA Numbers and Sets, Michaelmas 2025.
\newblock Example Sheet~1.
\newblock Copia consultada en \url{https://tartarus.org/gareth/maths/Numbers_and_Sets/NS_2025-26_q1.pdf}

\bibitem{cs103card}
Stanford University.
\newblock CS103: Mathematical Foundations of Computing, curso 2019--20.
\newblock Transparencias «Cardinality».
\newblock \url{https://web.stanford.edu/class/archive/cs/cs103/cs103.1204/}

\bibitem{cs103ps3}
Stanford University.
\newblock CS103: Mathematical Foundations of Computing, otoño de 2023.
\newblock Problem Set~3.
\newblock \url{https://web.stanford.edu/class/archive/cs/cs103/cs103.1242/psets/ps3}

\bibitem{cs103set}
Stanford University.
\newblock CS103: Mathematical Foundations of Computing, otoño de 2026.
\newblock Clase~8, «Set Theory Revisited».
\newblock \url{https://web.stanford.edu/class/cs103/}

\bibitem{hammack}
R.~Hammack.
\newblock \emph{Book of Proof}, 3.ª ed. (edición 3.4).
\newblock Publicación del autor, 2018.
\newblock \url{https://richardhammack.github.io/BookOfProof/}

\bibitem{mcs}
E.~Lehman, F.~T. Leighton y A.~R. Meyer.
\newblock \emph{Mathematics for Computer Science}.
\newblock Revisión del 18 de mayo de 2015. Texto del curso 6.042J del MIT.
\newblock \url{https://courses.csail.mit.edu/6.042/spring15/mcs.pdf}

\bibitem{mit6042}
A.~R. Meyer y A.~Chlipala.
\newblock 6.042J: Mathematics for Computer Science.
\newblock Massachusetts Institute of Technology, MIT OpenCourseWare, primavera de 2015.
\newblock \url{https://ocw.mit.edu/courses/6-042j-mathematics-for-computer-science-spring-2015/}

\bibitem{uw300}
University of Washington, Department of Mathematics.
\newblock Math 300: Introduction to Mathematical Reasoning, primavera de 2013.
\newblock Hoja «Problems on functions».
\newblock \url{https://sites.math.washington.edu/~conroy/2013/m300-spr2013/functions.pdf}

\bibitem{velleman}
D.~J. Velleman.
\newblock \emph{How to Prove It: A Structured Approach}, 3.ª ed.
\newblock Cambridge University Press, 2019.

\end{thebibliography}


\end{document}
